\documentclass[conference]{IEEEtran}
\IEEEoverridecommandlockouts

\usepackage{cite}
\usepackage{amsmath,amssymb}
\usepackage{booktabs}
\usepackage{graphicx}
\usepackage{url}
\usepackage{xcolor}

\title{Explainable Hybrid Detection of AI-Generated and Human-Written Content Under Leakage-Aware Evaluation}

\author{
\IEEEauthorblockN{Salma Hesham Salem}
\IEEEauthorblockA{
Natural Language Processing Project\\
Global AI vs. Human Content Detection 2026\\
}
}

\begin{document}
\maketitle

\begin{abstract}
The growth of generative models such as GPT, Gemini, and Claude has increased the need for reliable detectors that distinguish human-written content from AI-generated text and code. This paper presents an explainable hybrid detection framework for the Global AI vs. Human Content 2026 dataset and extends it with a leakage-aware comparative study. The first stage develops a DistilBERT--stylometric hybrid model that combines contextual transformer embeddings with interpretable machine-signature features such as lexical diversity, punctuation behavior, repetition, sentence length, and code-symbol usage. This model achieved 75.43\% accuracy, 0.754 macro-F1, and 0.879 ROC-AUC in the earlier cleaned random-split experiment. The second stage audits the benchmark more aggressively and shows that raw evaluation can be inflated by explicit ``AI-generated'' markers, duplicated templates, conflicting normalized texts, and metadata proxies such as source and AI model name. We therefore compare TF-IDF, SVM, stylometric-only, classical hybrid, frozen transformer, and transformer-stylometric ablations under random, prompt-grouped, and challenge-test settings. The combined study shows that hybrid semantic-stylometric modeling is promising, but also that AI-authorship detection results are only meaningful when leakage controls, ablation studies, and out-of-distribution checks are reported together.
\end{abstract}

\begin{IEEEkeywords}
AI-generated text detection, DistilBERT, stylometry, leakage audit, ablation study, machine signatures
\end{IEEEkeywords}

\section{Introduction}
Large language models such as GPT-4, Gemini, and Claude can produce fluent text that is difficult to distinguish from human-authored writing. This motivates binary AI-content detection systems for education, journalism, code review, and online trust. A robust detector should avoid relying on superficial dataset artifacts and should generalize across prompts, topics, and writing styles.

The project addresses the Global AI vs. Human Content 2026 task as a supervised binary classification problem. The postgrad focus is stylometric analysis: rather than only learning topic words, the detector should capture machine-signature patterns such as repetition, punctuation regularity, sentence-length smoothness, and lexical diversity.

This work is organized as a two-stage study. Stage 1 is the earlier hybrid modeling work, where DistilBERT embeddings were fused with stylometric features and achieved strong balanced results after basic leakage marker removal. Stage 2 is a stricter reproduction and audit, where we deliberately stress-test the benchmark using normalized de-duplication, metadata proxy analysis, prompt-grouped evaluation, sanity checks, and challenge examples. This makes the project stronger than a single high-score report because it explains both what works and when reported performance should not be trusted.

\section{Related Work}
Transformer encoders such as BERT~\cite{devlin2019bert}, based on the attention architecture~\cite{vaswani2017attention}, are strong baselines for text classification. Prior work on neural text authorship attribution shows that generated text can carry model-specific artifacts~\cite{uchendu2020authorship}. Statistical detectors such as GLTR~\cite{gehrmann2019gltr} study token likelihood patterns as evidence of machine generation. Broader analyses of language-model release risks also emphasize the difficulty of reliable detection under distribution shift~\cite{solaiman2019release,pan2024llmdetect}.

This work follows the same direction but adds a stronger comparative and ablation design. Instead of reporting only the final hybrid score, it separates the contribution of stylometric features, transformer embeddings, TF-IDF lexical features, and data-splitting assumptions.

\section{Datasets and Leakage Audit}
The raw dataset contains 20,000 rows with columns including content, label, type, source, prompt, topic, AI model, language, word count, character count, complexity score, and multiline-code flags. A direct audit found three major leakage risks.

First, every AI row contained an explicit ``AI-generated'' marker in the raw content. Second, the \texttt{source} and \texttt{ai\_model} columns were perfectly correlated with the target label. Third, normalized duplicate and conflicting examples were common. These findings make a naive random split unsuitable for publication-quality claims.

\begin{table}[htbp]
\caption{Cleaning Audit Summary}
\centering
\begin{tabular}{lr}
\toprule
Audit Item & Count \\
\midrule
Raw rows & 20000 \\
Rows with AI marker & 10004 \\
Conflicting normalized texts & 918 \\
Duplicate normalized rows before drop & 1121 \\
Rows after strict cleaning & 6406 \\
Rows removed total & 13594 \\
\bottomrule
\end{tabular}
\label{tab:cleaning}
\end{table}

After cleaning, the retained dataset was balanced, with 3,157 human and 3,249 AI samples. The strict cleaned subset contained only text-domain examples; the code subset was removed because normalized code templates frequently conflicted across labels. Therefore, this paper does not claim reliable code-domain generalization from the cleaned local data.

To strengthen the study, we additionally use RAID, a large robustness benchmark for AI-generated text detection. RAID is valuable because it covers multiple generators, domains, attacks, decoding settings, and repetition penalties. Since the full RAID benchmark contains millions of generations, our local experiment uses a reproducible streaming subset with 1,000 balanced abstract samples: 500 human and 500 AI examples. The AI subset contains 300 \texttt{llama-chat} examples and 200 \texttt{mpt} examples. This subset is not meant to replace the full benchmark; it is used to add a stronger external robustness experiment under local compute limits.

\section{Methodology}
\subsection{Preprocessing}
Preprocessing is intentionally light for modeling and stricter for auditing. For model input, explicit AI markers, URLs, email-like patterns, excessive whitespace, and very short samples are removed while preserving punctuation, capitalization, line breaks, and token structure. For audit-only de-duplication, text is normalized to detect repeated or conflicting templates. Labels are encoded as human = 0 and AI = 1. Metadata fields such as \texttt{id}, \texttt{source}, \texttt{prompt}, \texttt{ai\_model}, \texttt{word\_count}, \texttt{char\_count}, and language tags are excluded from model input.

\subsection{Stylometric Features}
The stylometric extractor computes character count, word count, sentence count, line count, average and standard deviation of word length, average and standard deviation of sentence length, lexical diversity, hapax ratio, word entropy, punctuation ratios, uppercase ratio, whitespace ratio, quote count, code-like brace/operator counts, repeated word ratio, repeated bigram ratio, repeated trigram ratio, and line-length statistics.

\subsection{Proposed DistilBERT--Stylometric Hybrid}
The proposed model combines semantic and stylometric views of the same input. First, each sample is tokenized using a DistilBERT tokenizer and passed through DistilBERT. Mean pooling is applied over token embeddings using the attention mask. Second, the stylometric vector is standardized and passed through a small projection layer. The pooled transformer representation and projected stylometric representation are concatenated and sent to a neural classification head. The model is trained using AdamW with learning rate $2\times10^{-5}$ and binary cross-entropy with logits loss. The best checkpoint is selected using validation macro-F1, and the decision threshold is tuned on validation predictions to avoid over-predicting one class.

\subsection{Comparative Models and Ablations}
The comparative study includes:
\begin{itemize}
\item TF-IDF word n-grams with Logistic Regression.
\item TF-IDF character n-grams with Logistic Regression.
\item TF-IDF word n-grams with calibrated Linear SVM.
\item Stylometric-only Logistic Regression.
\item Hybrid TF-IDF plus stylometric Logistic Regression.
\item Frozen DistilRoBERTa embeddings with Logistic Regression.
\item Frozen DistilRoBERTa embeddings plus stylometric features.
\item Fine-tuned DistilBERT plus stylometric neural fusion from the prior hybrid study.
\end{itemize}

The ablation logic is direct: stylometric-only tests whether machine signatures alone are separable; frozen transformer-only tests semantic representation without task fine-tuning; transformer plus stylometry tests feature fusion; and the final fine-tuned hybrid tests whether supervised transformer adaptation is necessary.

\section{Experimental Design}
Two protocols are reported. The earlier hybrid experiment used 12,800 training samples, 3,200 validation samples, and 4,000 test samples. This protocol evaluates the proposed full DistilBERT--stylometric neural model after removing explicit markers and excluding metadata from the model input.

The stricter leakage-aware pipeline uses a cleaned subset after normalized duplicate and conflict removal. It reports a stratified random split, a prompt-grouped split where prompt templates do not cross train and test partitions, and a separate challenge set containing unseen topics, messy human prose, structured AI prose, short examples, long examples, and code/text cases.

For RAID, three evaluation views are used: a random stratified split, a source-grouped split using source identifiers, and a model-holdout split where \texttt{llama-chat} generations are reserved for testing. The model-holdout split directly tests whether the detector generalizes to an unseen generator style.

We also perform a cross-dataset external benchmark: models trained on the original cleaned Global AI vs Human dataset are evaluated directly on the RAID subset without retraining. This tests whether the earlier detector family learned transferable AI-writing signals or dataset-specific artifacts.

Three sanity checks are used: a label-shuffle text model, a length-only model, and a prompt-only model. A scientifically valid setup should push these controls near chance.

\section{Results}
\subsection{Prior Hybrid Model Performance}
The earlier DistilBERT--stylometric hybrid selected epoch 2 based on validation macro-F1. Validation accuracy was 74.56\%, validation balanced accuracy was 74.56\%, validation macro-F1 was 0.7449, and validation ROC-AUC was 0.8711. The final test results are shown in Table~\ref{tab:prior-hybrid}. The confusion matrix in Table~\ref{tab:confusion} shows balanced behavior: the model does not collapse into predicting only AI or only human.

\begin{table}[htbp]
\caption{Prior Fine-Tuned DistilBERT--Stylometric Hybrid}
\centering
\begin{tabular}{lr}
\toprule
Metric & Value \\
\midrule
Accuracy & 75.43\% \\
Balanced accuracy & 75.42\% \\
AI F1-score & 0.764 \\
Macro-F1 & 0.754 \\
ROC-AUC & 0.879 \\
\bottomrule
\end{tabular}
\label{tab:prior-hybrid}
\end{table}

\begin{table}[htbp]
\caption{Prior Hybrid Confusion Matrix}
\centering
\begin{tabular}{lrr}
\toprule
 & Predicted Human & Predicted AI \\
\midrule
Actual Human & 1427 & 572 \\
Actual AI & 411 & 1590 \\
\bottomrule
\end{tabular}
\label{tab:confusion}
\end{table}

\subsection{Leakage-Aware Comparative Study}
Table~\ref{tab:main-results} reports the stricter reproduction after normalized duplicate/conflict removal. Scores are intentionally lower than the prior hybrid result because the stricter protocol removes many repeated templates and uses more conservative evaluation.

\begin{table*}[htbp]
\caption{Main Model Results After Strict Cleaning}
\centering
\begin{tabular}{llrrrr}
\toprule
Split & Model & Accuracy & Precision & Recall & F1 \\
\midrule
Random & TF-IDF Word LR & 0.517 & 0.524 & 0.513 & 0.519 \\
Random & TF-IDF Char LR & 0.495 & 0.502 & 0.493 & 0.497 \\
Random & TF-IDF Word SVM & 0.489 & 0.497 & 0.764 & 0.602 \\
Random & Stylometric LR & 0.512 & 0.518 & 0.540 & 0.529 \\
Random & Hybrid TF-IDF + Stylometric LR & 0.541 & 0.549 & 0.534 & 0.541 \\
Prompt-grouped & TF-IDF Word LR & 0.512 & 0.526 & 0.598 & 0.560 \\
Prompt-grouped & TF-IDF Char LR & 0.530 & 0.533 & 0.758 & 0.626 \\
Prompt-grouped & TF-IDF Word SVM & 0.509 & 0.517 & 0.825 & 0.636 \\
Prompt-grouped & Stylometric LR & 0.521 & 0.538 & 0.545 & 0.542 \\
Prompt-grouped & Hybrid TF-IDF + Stylometric LR & 0.504 & 0.520 & 0.587 & 0.551 \\
\bottomrule
\end{tabular}
\label{tab:main-results}
\end{table*}

\begin{table*}[htbp]
\caption{Ablation and Protocol Comparison}
\centering
\begin{tabular}{llrrr}
\toprule
Protocol & Model & Accuracy & F1 / Macro-F1 & ROC-AUC \\
\midrule
Prior cleaned random split & Fine-tuned DistilBERT + stylometry & 0.754 & 0.754 & 0.879 \\
Strict random split & Stylometric LR & 0.512 & 0.529 & 0.515 \\
Strict random split & Frozen DistilRoBERTa Embedding LR & 0.512 & 0.521 & 0.517 \\
Strict random split & Frozen DistilRoBERTa + stylometry & 0.507 & 0.520 & 0.517 \\
Strict random split & Hybrid TF-IDF + stylometry & 0.541 & 0.541 & 0.539 \\
Strict prompt-grouped split & TF-IDF Char LR & 0.530 & 0.626 & 0.523 \\
Strict prompt-grouped split & Frozen DistilRoBERTa + stylometry & 0.521 & 0.550 & 0.508 \\
\bottomrule
\end{tabular}
\label{tab:ablation}
\end{table*}

\begin{table}[htbp]
\caption{Challenge Set Results}
\centering
\begin{tabular}{lrr}
\toprule
Model & Accuracy & F1 \\
\midrule
TF-IDF Word LR & 0.750 & 0.783 \\
Stylometric LR & 0.500 & 0.500 \\
Hybrid TF-IDF + Stylometric LR & 0.300 & 0.417 \\
TF-IDF Char LR & 0.500 & 0.286 \\
TF-IDF Word SVM & 0.300 & 0.125 \\
\bottomrule
\end{tabular}
\label{tab:challenge}
\end{table}

\begin{table}[htbp]
\caption{Frozen DistilRoBERTa Embedding Results}
\centering
\begin{tabular}{lrr}
\toprule
Split and Model & Accuracy & F1 \\
\midrule
Random Embedding LR & 0.512 & 0.521 \\
Random Embedding + Stylometry & 0.507 & 0.520 \\
Prompt-grouped Embedding LR & 0.508 & 0.537 \\
Prompt-grouped Embedding + Stylometry & 0.521 & 0.550 \\
Challenge Embedding LR & 0.300 & 0.125 \\
\bottomrule
\end{tabular}
\label{tab:transformer}
\end{table}

\begin{table*}[htbp]
\caption{RAID Quick-Subset Robustness Results}
\centering
\begin{tabular}{llrrrr}
\toprule
Split & Model & Accuracy & Precision & Recall & Macro-F1 \\
\midrule
Random & RAID TF-IDF Word LR & 0.933 & 0.945 & 0.920 & 0.933 \\
Random & RAID TF-IDF Char LR & 0.880 & 0.861 & 0.907 & 0.880 \\
Random & RAID TF-IDF Word SVM & 0.980 & 0.986 & 0.973 & 0.980 \\
Source-grouped & RAID TF-IDF Word LR & 0.924 & 0.965 & 0.894 & 0.923 \\
Source-grouped & RAID TF-IDF Char LR & 0.901 & 0.939 & 0.878 & 0.901 \\
Source-grouped & RAID TF-IDF Word SVM & 0.964 & 0.991 & 0.943 & 0.964 \\
Held-out llama-chat & RAID TF-IDF Word LR & 0.691 & 1.000 & 0.177 & 0.551 \\
Held-out llama-chat & RAID TF-IDF Char LR & 0.706 & 0.922 & 0.237 & 0.592 \\
Held-out llama-chat & RAID TF-IDF Word SVM & 0.684 & 1.000 & 0.157 & 0.534 \\
\bottomrule
\end{tabular}
\label{tab:raid}
\end{table*}

\begin{table*}[htbp]
\caption{Previous Project Models Evaluated Directly on RAID}
\centering
\begin{tabular}{lrrrr}
\toprule
Detector & Accuracy & AI Recall & AI F1 & Macro-F1 \\
\midrule
TF-IDF Word LR & 0.463 & 0.924 & 0.632 & 0.318 \\
TF-IDF Char LR & 0.461 & 0.174 & 0.244 & 0.413 \\
TF-IDF Word SVM & 0.487 & 0.974 & 0.655 & 0.328 \\
Stylometric LR & 0.500 & 0.000 & 0.000 & 0.333 \\
Hybrid TF-IDF + Stylometric LR & 0.414 & 0.074 & 0.112 & 0.337 \\
Saved Epoch TF-IDF SGD & 0.454 & 0.898 & 0.622 & 0.320 \\
\bottomrule
\end{tabular}
\label{tab:previous-on-raid}
\end{table*}

\begin{table}[htbp]
\caption{Sanity Checks}
\centering
\begin{tabular}{lrr}
\toprule
Check & Accuracy & F1 \\
\midrule
Label-shuffle text model & 0.515 & 0.515 \\
Length-only model & 0.505 & 0.510 \\
Prompt-only model & 0.496 & 0.522 \\
\bottomrule
\end{tabular}
\label{tab:sanity}
\end{table}

\section{Discussion}
The central result is not simply that one model has the highest score. The stronger finding is that model quality and benchmark quality interact. The prior fine-tuned DistilBERT--stylometric hybrid achieved the best balanced performance, supporting the project hypothesis that semantic representations and machine-signature features complement each other. However, the stricter audit shows that a dataset can become much harder once explicit markers, metadata proxies, duplicated templates, and conflicting normalized texts are removed.

The ablation study clarifies the role of each component. Stylometric-only classification is useful for explainability but weak by itself. Frozen transformer embeddings also remain weak, showing that generic representations without task-specific fine-tuning are not enough for this dataset. Classical TF-IDF models remain competitive, especially under prompt-grouped evaluation, which suggests that lexical and phrase-level artifacts still matter. The full fine-tuned hybrid is therefore the most defensible proposed model, but its results should be reported together with leakage controls and stricter validation.

The RAID experiment improves the paper's comparative study. On the random RAID split, TF-IDF Word SVM reaches 0.980 macro-F1. On the source-grouped split, it remains high at 0.964 macro-F1. However, when \texttt{llama-chat} is held out as an unseen generator, macro-F1 drops to 0.534--0.592 depending on the baseline. This gap is important: it shows that high random-split performance does not imply robust cross-generator detection. The model-holdout setting is therefore the most meaningful RAID result for discussion.

The cross-dataset experiment is even more important for judging the previous project models. When models trained on the original cleaned dataset are tested directly on RAID, the best macro-F1 is only 0.413. Some models achieve high AI recall but very poor macro-F1, meaning they mostly over-predict AI and fail to distinguish both classes reliably. This confirms that the earlier dataset does not teach robust detector behavior on its own. RAID should therefore be used as the main benchmark for final claims, while the original dataset should be framed as a leakage case study and development set.

The practical conclusion is that the system should be presented as an explainable research detector, not as a guaranteed real-world AI-authorship judge. The Streamlit platform demonstrates the detector, exposes stylometric signals, and shows the training log/loss curve so users can inspect model behavior instead of seeing only a final label.

\section{Limitations}
The strict cleaned subset loses the code domain because many code samples are duplicated or label-conflicting after normalization. The challenge set is useful for stress testing but is small. The RAID run reported here is a quick local subset, not the full 7.4M-row benchmark. It currently covers the abstracts domain and two AI generators, so it should be expanded to more RAID domains, attacks, and held-out models before final publication. The earlier full fine-tuned hybrid and the later strict reproduction are not identical protocols, so their comparison should be interpreted as a protocol comparison rather than a single controlled leaderboard.

\section{Conclusion}
This paper presents an explainable hybrid approach for detecting AI-generated versus human-written content and strengthens it with a leakage-aware comparative study. The prior fine-tuned DistilBERT--stylometric model achieved 75.43\% accuracy, 0.754 macro-F1, and 0.879 ROC-AUC, showing that semantic-stylometric fusion is promising. The stricter reproduction and RAID experiments show why such results must be accompanied by leakage audits, ablations, grouped splits, sanity checks, challenge tests, and generator-holdout evaluation. The final contribution is therefore both a proposed hybrid detector and a defensible evaluation protocol for AI-authorship detection.

\bibliographystyle{IEEEtran}
\bibliography{references}

\end{document}
